\section{Related Work}
\label{sec:related}

Table~\ref{tab:comparison} summarizes the structural differences across contemporary secret-handling and agent-isolation paradigms. While existing defenses tackle individual aspects of the problem---such as post-hoc taint tracking, heuristic prompt rewriting, or planner-level tool mediation---none simultaneously guarantee zero plaintext observation and destination-bound resolution for vision-driven GUI agents.

\paragraph{Placeholder and anonymization defenses.} Mobile placeholder schemes replace sensitive text with type-preserving tokens and reverse the mapping at entry \cite{zhao2026anonymization}; recent work resolves placeholders locally but is detector-based, session-scoped, and validates index ranges rather than target destinations. We differ fundamentally on mint source (authoritative OS source, not screen scraping), unforgeability, destination-and-task binding, typed lifecycle, and a threat model that includes redirection. Unbound placeholders remain completely vulnerable to Confused Deputy redirection ($30/30$ leakage under Attack~2), whereas \SH{} prevents redirection at the OS boundary.

\paragraph{Capability and isolation approaches.} Planner-side capability schemes and dual-model patterns hide values from planners in tool-calling settings \cite{debenedetti2025camel,costa2025fides}; isolation designs restrict tool permissions \cite{bagdasarian2024airgap,wu2025isolategpt,shi2025progent}. These operate in planning or tool layers rather than where a secret enters a UI field, assuming structured APIs rather than graphical mobile environments where untrusted vision agents observe raw pixels and actuate via accessibility.

\paragraph{Prompt sanitizers.} Detection-and-rewrite schemes suppress sensitive spans in prompts \cite{siyan2025papillon,zeng2025privacyrestore,chong2024casper}. They reduce model-side exposure but lack destination binding and OS enforcement. Furthermore, prompt sanitization is probabilistic: adversarial formatting and obfuscation can evade detection filters.

\paragraph{Platform mechanisms and reference monitors.} Android's SMS Retriever/Consent \cite{androidSmsRetriever,androidSmsConsent}, sensitive-notification OTP pipeline \cite{android15privacy}, Credential Manager \cite{androidCredentialManager}, and Protected Confirmation \cite{androidProtectedConfirmation} are valuable primitives, but trust requesting apps without intra-app or post-entry confinement. \SH{} applies complete mediation and least privilege \cite{saltzer1975protection} to UI injection by untrusted agents, building on capability systems (Capsicum \cite{watson2010capsicum}, Dennis and Van Horn \cite{dennis1966semantics}), Android app isolation \cite{felt2011permission,bugiel2013flexible,rossi2021seapp}, and trusted UI paths (SchrodinText \cite{continella2020schrodintext}, Rushmore \cite{andrus2015rushmore}, GYRUS \cite{jang2014gyrus}). While Cloak \& Dagger \cite{fratantonio2017cloak} revealed accessibility overlay abuses, \SH{} leverages overlays to shield secrets from vision agents. Unlike post-hoc taint tracking (TaintDroid \cite{enck2010taintdroid}, DroidScope \cite{yan2012droidscope}) or endpoint token binding \cite{rfc8707,rfc9449,rfc9700,webauthn3,nist80063b,rfc4226,rfc6238}, \SH{} binds on-device resolution directly to verified UI destinations under documented agent prompt-injection threats \cite{du2025mobileattacks,liao2024eia,chen2025fineprint,greshake2023indirect,debenedetti2024agentdojo,zhan2024injecagent}.
